\section*{Problem 8: Decoding III --- Chien Search \& Magnitudes (Level 8)}

\subsubsection*{Mathematical part}
\textbf{Recall:} Once the locator polynomial $\Lambda(x)$ has been constructed, the remaining task is to find its roots. Instead of employing complex root-finding algorithms, we leverage the finite nature of our field.

\textbf{Chien Search:} Simply evaluate $\Lambda(x)$ at every non-zero element of the field. Whenever $\Lambda(a) = 0$, we have discovered an error location! Because the roots are $X_i^{-1}$, if $a$ is a root, the error position is given by the corresponding inverse power.

\textbf{Error Magnitudes:} Knowing \emph{where} an error occurred is only half the problem. We must also determine \emph{how much} the received value should be changed. Let $Y_i$ be the magnitude of the error at position $X_i$. The syndromes were defined as $S_k = r(\alpha^k)$. Because the transmitted codeword vanishes at these roots, the syndromes depend entirely on the errors:
\[
S_k = Y_1 X_1^k + Y_2 X_2^k + \dots + Y_v X_v^k
\]
Since we now know the error locations $X_i$, this is simply another system of linear equations (a Vandermonde matrix) where the only unknowns are the magnitudes $Y_i$. We can solve this using the exact same Gaussian elimination utility!

\paragraph{Worked Example}
Suppose Chien Search found $v=2$ errors at locations $X_1$ and $X_2$. The syndromes $S_0, S_1$ give us the following linear system for the unknown magnitudes $Y_1, Y_2$:
\[
\begin{bmatrix}
X_1^0 & X_2^0 \\
X_1^1 & X_2^1
\end{bmatrix}
\begin{bmatrix}
Y_1 \\
Y_2
\end{bmatrix}
=
\begin{bmatrix}
S_0 \\
S_1
\end{bmatrix}
\]
Which simplifies to:
\[
\begin{bmatrix}
1 & 1 \\
X_1 & X_2
\end{bmatrix}
\begin{bmatrix}
Y_1 \\
Y_2
\end{bmatrix}
=
\begin{bmatrix}
S_0 \\
S_1
\end{bmatrix}
\]
Notice that the matrix is purely built from powers of the roots $X_i$.

\subsubsection*{Implementation part}
\textbf{Task:} Implement \lstinline|chien_search| and \lstinline|linear_error_magnitudes|.

\textbf{Details:} 
\begin{itemize}
    \item \lstinline|chien_search|: Loop over every position $i$ in the message (from $0$ to $msg\_len - 1$). The root candidate is simply \lstinline|field.alpha ** -i|. Your \lstinline|__pow__| from Level~1 already handles negative exponents natively via Fermat's Little Theorem, so no manual conversion is needed. Evaluate the locator polynomial at this root. If it evaluates to zero (\lstinline|not err_loc(root)| using our duck-typed \lstinline|__bool__|), append $i$ to the list of error positions.
    \item \lstinline|linear_error_magnitudes|: If no errors were found ($t = \text{len(err\_pos)} == 0$), return an empty list \lstinline|[]|. Otherwise, we must solve the system $H' Y = \vec{S}'$ to find the error magnitudes $Y_i$:
\[
\underbrace{
\begin{bmatrix}
X_1^0 & X_2^0 & \cdots & X_t^0 \\
X_1^1 & X_2^1 & \cdots & X_t^1 \\
\vdots & \vdots & \ddots & \vdots \\
X_1^{t-1} & X_2^{t-1} & \cdots & X_t^{t-1}
\end{bmatrix}
}_{H'}
\begin{bmatrix}
Y_1 \\
Y_2 \\
\vdots \\
Y_t
\end{bmatrix}
=
\underbrace{
\begin{bmatrix}
S_0 \\
S_1 \\
\vdots \\
S_{t-1}
\end{bmatrix}
}_{\vec{S}'}
\]
Here, $X_i = \alpha^{pos_i}$ for each position in \lstinline|err_pos|. Notice that the vector $\vec{S}'$ is simply the first $t$ syndromes. Build your matrix $H'$ and vector $\vec{S}'$, call \lstinline|solve_linear| on them to get the magnitudes $Y$, and simply return the list $Y$.
\end{itemize}

\begin{center}
\renewcommand{\arraystretch}{1.5}
\begin{tabularx}{\textwidth}{@{} >{\bfseries}l X c @{}}
\toprule
Function & Arguments & Returns \\ \midrule
chien\_search & \lstinline|err_loc: Polynomial, msg_len: int, field: ExtensionField| & \lstinline|list[int]| \\
linear\_error\_magnitudes & \lstinline|syndromes: list[GaloisFieldElement], err_pos: list[int], field: ExtensionField| & \lstinline|list[GaloisFieldElement]| \\
\bottomrule
\end{tabularx}
\end{center}

\subsubsection*{Experimentation part}
\textbf{UI Guide:} Open the dashboard and watch the final magic. By solving the second linear system, the UI will correctly report exactly how much to subtract from your corrupted bytes to perfectly restore the original text message.

\vspace{1cm}
